\section{Introduction}\label{sec:intro}

Large language model agents act for their users by
looping through observation, reasoning, and action in an
external environment.
What such an agent believes about the information it
reads enters tool calls and operations directly, which
makes the reliable interpretation of external content a
precondition for execution safety.
Prior work has shown that indirect prompt injections
buried in web pages and tool returns can steer an agent
away from the user's intent \citep{zhan2024injecagent,
evtimov2025wasp, johnson2025manipulating}, that
perturbations of visual input can hijack multi-step
web agents \citep{wu2024dissecting}, and that deceptive
interface design can be reproduced and probed
systematically \citep{wang2025agentvigil}.
These findings indicate that an environment does not only
intervene on an agent through explicit instructions; it
can also act on the way the agent understands and uses
information, and thereby shape the decisions the agent
makes autonomously.

Charts are a prominent part of such external
information.
They organize data into visual forms that support
comparison and inference, and they lend an appearance of
evidential support to whatever judgment is attached to
them.
Yet the basis for a correct reading and a misleading
presentation can coexist in the same figure: research in
visualization shows that charts are used to support
misleading inferences even when the data are drawn
faithfully \citep{lisnic2023misleading, nadib2026guardrail}.
A curve that moves upward across the rendered image is
observed accurately, for instance, but whether it
represents growth depends on the axis mapping currently
in force; the meaning of color intensity and bar height
is likewise read through the encoding relations of the
present figure.
A misleading presentation can induce an agent to apply
its customary reading rules unchanged, converting
genuinely visible evidence into a wrong conclusion.
\citet{chen2025unmasking} document that multimodal
models are systematically vulnerable to such
manipulations in chart question answering,
\citep{mahbub2025perils, mahbub2026chart} trace the
vulnerability and its partial mitigation across
model families, and our companion work under review
formalizes the execution-level consequence: on a paired
benchmark of executable web tasks, biases induced by
misleading charts propagate into subsequent actions and
change task outcomes.\footnote{Companion submission
under review; the paired benchmark \bench{} and the
execution-level evaluation protocol it introduces are
described where first used in this paper.}

Against misleading charts, model-level defenses have
improved chart question answering through tabular
answering and re-plotting procedures that force a
structured restatement of the evidence
\citep{tonglet2025protecting}.
These procedures provide an effective basis for
model-level repair, but the reliability of the
verification and correction itself remains open.
In the inference-only setting, intrinsic
self-correction can damage answers that were initially
correct \citep{huang2023large}, and when the evaluation
component is itself exposed to the environment,
reflection and search can introduce new attack surfaces
rather than remove existing ones \citep{wu2024dissecting}.
Chart-domain studies observe a further gap: a model may
correctly identify chart properties during extraction
and still fail to let those properties govern the
downstream inference \citep{mahbub2026chart}.
These observations point to two connected requirements
for a defense: correction must rest on evidence that can
be checked, and a corrected interpretation must
actually change subsequent judgments.
Repeated observation or an additional review round can
discover an error, but it can equally re-adopt an
inapplicable reading, or overwrite a judgment that was
already correct.

Stateful execution raises requirements beyond
chart question answering: an understanding formed from
one chart feeds later planning and operations even
after the chart has left the viewport.
Once an axis mapping has been verified, later judgments
should keep using that verification; when sufficient
new evidence appears, the same interpretation should
become revisable.
Reflexion-style memory carries verbalized feedback into
later attempts \citep{shinn2023reflexion}, and recent
agent-memory systems organize experience into
retrievable state \citep{wang2024survey, hu2025agents};
an inductive chart requires something more specific,
the current validity of interpretations: which rules
have evidential support, which have been refuted, and
what new evidence suffices to change an accepted
reading?
We pose one unified question: \emph{how can an agent
reliably revise visual interpretation rules, and keep
validated rules in force over subsequent execution?}

We answer this question with a defense that combines
\emph{competing explanation chains} with a
\emph{persistent rule state}.
When the agent forms a chart-dependent judgment, the
harness expresses it as $O\land B\Rightarrow C_q(a)$:
an observation $O$ in the chart, an interpretation rule
$B$ that turns the observation into task meaning, and
the conclusion $C_q(a)$ that an action $a$ requires
under the task $q$.
The harness then builds candidate chains with different
interpretive bases on the same complete chart, turning
default visual conventions into explicit objects of
test, and compares their applicability using evidence
drawn from the figure itself.
Validated rule revisions enter the agent's runtime
harness as an interpretation state that later reasoning
reads and updates by scope; subsequent decisions reuse
the rules currently in force, and re-enter verification
only when new effective evidence challenges them.
Competing explanations determine why an interpretation
needs to change; the persistent state carries that
evidenced change into later actions.

The contributions of this paper are two-fold.
First, a correction paradigm based on competing
explanation chains: the relation
$O\land B\Rightarrow C_q(a)$ unifies visual observation,
interpretation rules, and action judgments; competing
chains are constructed on one chart, and the
applicability of customary readings is tested against
in-chart evidence, so that the rules responsible for
wrong inferences are identified and revised.
Second, a stateful defense for continuing execution:
validated rules, their revision evidence, and their
scope conditions are maintained as task-scoped state
inside the harness, so that later decisions reuse
confirmed chart readings and revise them only on new
effective evidence.
